\documentclass[11pt]{article}

\usepackage[margin=1.1in]{geometry}
\usepackage{amsmath}
\usepackage{amssymb}
\usepackage{amsthm}
\usepackage{mathrsfs}
\usepackage{booktabs}
\usepackage{enumitem}
\usepackage{tikz}
\usepackage[hidelinks]{hyperref}
\usepackage[activate={true,nocompatibility},expansion=false]{microtype}

\theoremstyle{plain}
\newtheorem{theorem}{Theorem}[section]
\newtheorem{proposition}[theorem]{Proposition}
\newtheorem{lemma}[theorem]{Lemma}
\newtheorem{corollary}[theorem]{Corollary}
\newtheorem{conjecture}[theorem]{Conjecture}
\theoremstyle{definition}
\newtheorem{definition}[theorem]{Definition}
\newtheorem{example}[theorem]{Example}
\newtheorem{decision}[theorem]{Decision}
\theoremstyle{remark}
\newtheorem{remark}[theorem]{Remark}

% ---- source calculus, following the companion note ----
\newcommand{\quo}[1]{@#1}
\newcommand{\drp}[1]{{*}#1}
\newcommand{\inp}[3]{\mathsf{for}(\,#1 \leftarrow #2\,)#3}
\newcommand{\out}[2]{#1!(#2)}
\newcommand{\nil}{\mathbf{0}}
\newcommand{\enc}[1]{[\![#1]\!]}
\newcommand{\nequiv}{\equiv_{\mathcal{N}}}
\newcommand{\Ctx}{\mathscr{C}}
\newcommand{\Obs}{\mathcal{D}}
\newcommand{\bisim}{\approx}
\newcommand{\nms}[1]{\mathcal{N}^{*}(#1)}

% ---- combinators ----
\newcommand{\mm}{\mathsf{m}}
\newcommand{\dd}{\mathsf{d}}
\newcommand{\kk}{\mathsf{k}}
\newcommand{\fw}{\mathsf{fw}}
\newcommand{\bl}{\mathsf{b}_{\mathsf{l}}}
\newcommand{\br}{\mathsf{b}_{\mathsf{r}}}
\newcommand{\sy}{\mathsf{s}}
\newcommand{\ev}{\mathsf{e}}
\newcommand{\qq}{\mathsf{q}}
\newcommand{\cons}{\mathsf{cons}}
\newcommand{\consp}{\cons_{\mid}}
\newcommand{\consm}{\cons_{\mm}}
\newcommand{\consd}{\cons_{\dd}}
\newcommand{\conss}{\cons_{\sy}}
\newcommand{\red}{\longrightarrow}
\newcommand{\wred}{\longrightarrow^{*}}
\newcommand{\scong}{\equiv}

% ---- logic ----
\newcommand{\env}{\mathcal{E}}
\newcommand{\Frm}{\mathcal{F}}
\newcommand{\Nmf}{\mathcal{V}}
\newcommand{\may}[1]{\langle #1 \rangle}
\newcommand{\mst}[1]{[ #1 ]}
\newcommand{\rcp}[1]{(\!|\, #1 \,|\!)}
\newcommand{\sat}{\models}
\newcommand{\chf}[1]{\chi(#1)}
\newcommand{\Occ}{\mathsf{Occ}}
\newcommand{\Conf}{\mathsf{Conf}}
\newcommand{\Cl}{\mathsf{Cl}}
\newcommand{\Atm}{\mathcal{A}}
\newcommand{\atoms}[1]{\mathrm{at}^{*}(#1)}
\newcommand{\comps}[1]{\mathrm{comps}(#1)}
\newcommand{\Sub}[1]{\mathrm{Sub}^{*}(#1)}
\newcommand{\shp}{\mathsf{NM}}
\newcommand{\fix}{\mu}
\newcommand{\sem}[1]{[\![#1]\!]}

\title{\textbf{A Spatial-Behavioural Logic for the Rho Combinators}\\[4pt]
\large Structural formulae, one context modality, and decidable occupancy
of the expressiveness lattice}
\author{F1R3FLY.io research note}
\date{September 2026 --- draft 1}

\begin{document}
\maketitle

\begin{abstract}
\noindent
The companion note \cite{rhocomb3} repairs the name-free concurrent combinators
for the rho calculus and replaces the separation theorem with an expressiveness
lattice over subsets of the atom set. Each edge of that lattice is witnessed by
an invariant --- name closure, atom count, message count --- established in the
metatheory and stated in prose. This note supplies a logic in which those
invariants are formulae. The logic is the usual propositional connectives,
the spatial pair $\nil$ and $\mid$, one structural production per combinator
shape with name positions written $\quo{\phi}$, a single context modality
$\may{\phi}\psi$, and a parameterised least fixpoint over namespaces. It has no
restriction connectives, no guarantee, and no name sort beyond $\quo{\phi}$,
because the calculus has no binders and every name is a quotation.

Three things follow. Characteristic formulae exist and are exact, so the logic
separates terms precisely up to structural congruence, and nominals, counting
and the temporal operators are all derived rather than primitive
(\S\ref{sec:derived}). Occupancy of a lattice point --- the assertion that a
term is a term of the sub-calculus over an atom set $\Atm$ --- is the nullary
fixpoint $\Occ(\Atm)$, is checkable in one pass over the term
(Theorem~\ref{thm:occlinear}), and is invariant under reduction exactly when
$\Atm$ is rule-closed (Proposition~\ref{prop:subjred}). The constructor axis
has a second and sharper formula, the namespace closure $\Cl_{\Atm}(T)$, which
is the fixpoint whose disjuncts are in bijection with the constructors present;
it is the logical form of the name-closure lemma. Third, rule-closure has a
consequence for the lattice as drawn: at the eight points carrying $\consd$
without $\dd$, occupancy is not invariant, so the two axes of
\cite[Fig.~2]{rhocomb3} are not independent for this predicate
(Proposition~\ref{prop:notindep}). The modality-free fragment is decidable
outright; the full logic is decidable on the $\ev$-free sublattice modulo one
finite-witness obligation, and undecidable above it.
\end{abstract}

\tableofcontents

%==============================================================================
\section{What this note is}
\label{sec:what}

The expressiveness lattice of \cite[\S10]{rhocomb3} is a family of sub-calculi
of the rho combinators, ordered by inclusion of atom sets, with each strict edge
witnessed by a behaviour available above it and not below. The witnesses are
produced by three invariants: names reachable from a term lie in the closure of
its names under the node shapes of the constructors present; every rule but
$\ev$'s strictly decreases the count of non-message atoms; and $\dd$ is the only
free source of copies.

Those invariants are stated about the calculus. This note states them
\emph{in} a logic over the calculus --- so that ``$T$ sits at or below lattice
point $S$'' is a formula, and so that checking it is a decision procedure rather
than a proof. That is the whole of the brief. In particular:

\begin{itemize}[leftmargin=1.4em,itemsep=2pt,topsep=3pt]
\item This is a logic, not a type system. Nothing here is a typing judgement,
      nothing is syntax-directed on a derivation, and there is no subject-reduction
      theorem for terms-in-context. The intended route to a type system is
      generation from the logic, which is out of scope here.
\item The logic is deliberately intensional. Theorem~\ref{thm:char} says it
      separates terms exactly up to $\scong$, which is far finer than the
      equivalence $\bisim_{\enc{\Ctx_{\mathrm{rho}}}}$ of
      \cite[Thm.~7.5]{rhocomb3}. That is wanted for occupancy and is wrong for
      anything else; \S\ref{sec:notdo} says so plainly.
\item Occupancy formulae must be decidable. They are, and better than decidable:
      linear (Theorem~\ref{thm:occlinear}).
\end{itemize}

The logic is a near relative of the hybrid logic of the F1R3Space semantics note
\cite{f1r3space}, and diverges from it in four places, each traceable to a
feature of the target calculus rather than to taste: there is no restriction, so
revelation, hiding, freshness and $\mathsf{free}(a)$ all go; there is no
unforgeable name, so the name sort collapses into $\quo{(-)}$ and needs no
Boolean structure of its own (Proposition~\ref{prop:namecollapse}); the
guarantee $\phi \rhd_{\eta} \psi$ is absent, which is what keeps the structural
fragment decidable; and the three action modalities are replaced by one context
modality rather than by the four of \cite{softwhere}. The last is a restriction
and is discussed in \S\ref{sec:notdo}.

%==============================================================================
\section{The calculus recalled}
\label{sec:calc}

We follow \cite[\S3]{rhocomb3} exactly, and repeat only what the logic needs.

\begin{definition}[Syntax]
\label{def:syntax}
\begin{align*}
 A &\;::=\; \mm(a,b) \;\mid\; \dd(a,b,c) \;\mid\; \kk(a) \;\mid\; \fw(a,b)
       \;\mid\; \bl(a,b) \;\mid\; \br(a,b) \;\mid\; \sy(a,b,c)\\
   &\qquad \mid\; \ev(a) \;\mid\; \qq(a,P)\\
   &\qquad \mid\; \consp(a,b,c) \;\mid\; \consm(a,b,c)
       \;\mid\; \consd(a,b,c,e) \;\mid\; \conss(a,b,c,e)\\
 P,Q &\;::=\; \nil \;\mid\; A \;\mid\; P \mid Q
 \qquad\qquad
 a,b,c \;::=\; \quo{P}
\end{align*}
\end{definition}

Write $\Atm_{\mathrm{full}}$ for the set of thirteen atom shapes, and
$\mathrm{ar}(\alpha)$ for the arity of $\alpha$ as a term former. Every argument
of every shape is a name except the second argument of $\qq$, which is a process.

\begin{definition}[Structural congruence]
$\scong$ is the least congruence making $(\mathsf{Proc},\mid,\nil)$ an abelian
monoid. There are no binders, hence no $\alpha$-equivalence, and no clause
identifying a dropped quotation with its content.
\end{definition}

\begin{definition}[Components]
$\comps{P}$ is the multiset of atoms at the top level of $P$ after erasure of
$\nil$; it is a complete invariant of $\scong$, so $P \scong Q$ iff
$\comps{P} = \comps{Q}$ as multisets, with equality of atoms taken
componentwise up to $\nequiv$ on names and $\scong$ on the payload of $\qq$.
\end{definition}

\begin{proposition}[Name equivalence on the target is structural]
\label{prop:nequiv}
On combinator names, $\quo{P} \nequiv \quo{Q}$ iff $P \scong Q$.
\end{proposition}

\begin{proof}
The source clause $\quo{\drp{x}} \nequiv x$ has no instance here: $\drp{(-)}$ is
not a process former of Definition~\ref{def:syntax}, the drop having been
demoted to the atom $\ev$ \cite[Rem.~3.2]{rhocomb3}. What remains of the
definition of $\nequiv$ is the congruence generated by $\scong$ under
quotation.
\end{proof}

Proposition~\ref{prop:nequiv} discharges corrigendum~4 of
\cite[\S12]{rhocomb3}, which observes that $\nequiv$ is used in the side
condition on observation and never defined for the combinators. It also says
that quotation is injective and onto the name universe, which is what
\S\ref{sec:logic} trades on and what the source calculus does not enjoy.

\begin{definition}[Reduction]
\label{def:red}
\[
\begin{array}{rcl@{\qquad}rcl}
  \dd(a,b,c) \mid \mm(a,v) &\red& \mm(b,v) \mid \mm(c,v)
    & \br(a,b) \mid \mm(a,v) &\red& \fw(b,v)\\
  \kk(a) \mid \mm(a,v) &\red& \nil
    & \bl(a,b) \mid \mm(a,v) &\red& \fw(v,b)\\
  \fw(a,b) \mid \mm(a,v) &\red& \mm(b,v)
    & \sy(a,b,c) \mid \mm(a,v) &\red& \fw(b,c)\\
  \ev(a) \mid \mm(a,\quo{P}) &\red& P
    & \fw(a,b) \mid \qq(a,P) &\red& \mm(b,\quo{P})\\
  \multicolumn{6}{l}{\consp(a,b,c) \mid \mm(a,\quo{P}) \mid \mm(b,\quo{Q})
     \ \red\ \mm(c,\quo{(P \mid Q)})}\\
  \multicolumn{6}{l}{\consm(a,b,c) \mid \mm(a,u) \mid \mm(b,v)
     \ \red\ \mm(c,\quo{\mm(u,v)})}\\
  \multicolumn{6}{l}{\consd(a,b,c,e) \mid \mm(a,u) \mid \mm(b,v)
     \mid \mm(c,w) \ \red\ \mm(e,\quo{\dd(u,v,w)})}\\
  \multicolumn{6}{l}{\conss(a,b,c,e) \mid \mm(a,u) \mid \mm(b,v)
     \mid \mm(c,w) \ \red\ \mm(e,\quo{\sy(u,v,w)})}
\end{array}
\]
closed under $\mid$ and under $\scong$, with subjects matched up to $\nequiv$.
\end{definition}

Every rule has between two and four premises, each premise an atom, and the
premises of a rule share names. That is the only structural fact the modality of
\S\ref{sec:logic} uses.

%==============================================================================
\section{The logic}
\label{sec:logic}

\subsection{Grammar}

\begin{definition}[Formulae]
\label{def:grammar}
Two sorts. Name formulae (namespaces) $\nu$, and process formulae $\phi$:
\[
\begin{array}{rcll}
  \nu &::=& \quo{\phi} \;\mid\; U & \text{namespace, namespace parameter}\\[4pt]
  \phi,\psi &::=& \mathsf{true} \;\mid\; \neg\phi \;\mid\; \phi \wedge \psi
        \;\mid\; \phi \vee \psi \;\mid\; \phi \Rightarrow \psi
        & \text{propositional}\\
     &\mid& \nil \;\mid\; \phi \mid \psi & \text{spatial}\\
     &\mid& \mm(\nu_1,\nu_2) \;\mid\; \dd(\nu_1,\nu_2,\nu_3)
            \;\mid\; \kk(\nu_1) & \text{structural}\\
     &\mid& \fw(\nu_1,\nu_2) \;\mid\; \bl(\nu_1,\nu_2)
            \;\mid\; \br(\nu_1,\nu_2)&\\
     &\mid& \sy(\nu_1,\nu_2,\nu_3) \;\mid\; \ev(\nu_1)
            \;\mid\; \qq(\nu_1,\phi)&\\
     &\mid& \consp(\nu_1,\nu_2,\nu_3) \;\mid\; \consm(\nu_1,\nu_2,\nu_3)&\\
     &\mid& \consd(\nu_1,\nu_2,\nu_3,\nu_4)
            \;\mid\; \conss(\nu_1,\nu_2,\nu_3,\nu_4)&\\
     &\mid& \may{\phi}\psi & \text{context modality}\\
     &\mid& \fix X(\vec U).\,\phi \;\mid\; X(\vec\nu)
            & \text{fixpoint, application}
\end{array}
\]
In $\fix X(\vec U).\,\phi$ the parameters $\vec U$ are bound in $\phi$, and $X$
must occur positively in $\phi$; the application $X(\vec\nu)$ must supply as
many namespaces as $X$ was declared with.
\end{definition}

The structural productions are in bijection with the atom shapes of
Definition~\ref{def:syntax}, and each has the arity of its shape. The one
asymmetry is $\qq(\nu_1,\phi)$, whose second argument is a process formula. It
is not a defect of the grammar but a faithful record of the calculus: $\qq$ is
the introduction form for a quotation, and its payload has not yet been quoted.

\begin{remark}[The gate migrates a formula across the sorts]
\label{rem:migrate}
Under the release rule $\fw(a,b) \mid \qq(a,P) \red \mm(b,\quo{P})$ the payload
is described by $\psi$ before the step and by $\quo{\psi}$ after. Formally,
$\qq(\nu,\psi)$ is carried to $\mm(\nu',\quo{\psi})$ where $\nu'$ is the
forwarder's target. The logic is thereby able to watch the sort change that
quotation performs, and it is the only place in the grammar where that change is
visible. A grammar that wrote $\qq(\nu_1,\nu_2)$ would be uniform and would be
uninhabited: no well-formed $\qq$ term has a name in its second argument.
\end{remark}

\begin{decision}[No guarantee, no restriction, one modality]
\label{dec:omit}
$\phi \rhd_{\eta} \psi$ is omitted, and with it the quantification over an
infinite class of environments that makes spatial model checking hard; the
single-step consequence of it that this setting wants is recovered by
$\may{\phi}\psi$. The restriction connectives of \cite{f1r3space} are omitted
because the calculus has no $\mathsf{new}$. The modality set of
\cite{softwhere} is cut from four to one; \S\ref{sec:notdo} records the cost.
\end{decision}

\subsection{Satisfaction}

Satisfaction is defined on closed terms and closed formulae, relative to an
environment $\env$ interpreting free fixpoint variables. We suppress $\env$
except in the fixpoint clauses.

\begin{definition}[Name sort]
\label{def:satname}
\[
  a \sat \quo{\phi} \quad\text{iff}\quad
  a \nequiv \quo{P} \ \text{for some } P \text{ with } P \sat \phi .
\]
\end{definition}

\begin{proposition}[The name sort is definable]
\label{prop:namecollapse}
Quotation is a bijection from $\scong$-classes of processes to $\nequiv$-classes
of names, and $\sem{\quo{(-)}}$ preserves and reflects Boolean structure:
\[
  \quo{(\phi \wedge \psi)} \equiv \quo{\phi} \wedge \quo{\psi},
  \qquad
  \quo{(\phi \vee \psi)} \equiv \quo{\phi} \vee \quo{\psi},
  \qquad
  \quo{(\neg\phi)} \equiv \neg\quo{\phi},
\]
and $\quo{\mathsf{true}}$ is the namespace of all names.
\end{proposition}

\begin{proof}
Bijectivity is Proposition~\ref{prop:nequiv} together with the fact that every
name of Definition~\ref{def:syntax} is a quotation. The four equivalences follow
by unfolding Definition~\ref{def:satname} against that bijection.
\end{proof}

So the name sort carries no productions of its own beyond $\quo{(-)}$ and the
parameter $U$, and the Boolean namespace algebra that \cite{f1r3space} must add
to the name sort is here a theorem. The reason is exactly that the combinators
have no unforgeable names: in \cite{f1r3space} $a \sat \quo{\phi}$ is false at
every atom, which forces $\top_{\mathcal{N}}$ to be primitive and a separate
Boolean algebra with it.

\begin{definition}[Process sort]
\label{def:satproc}
Propositional clauses are standard. Otherwise:
\[
\begin{array}{lcl}
  P \sat \nil &\text{iff}& P \scong \nil\\[2pt]
  P \sat \phi \mid \psi &\text{iff}&
     \exists P_1,P_2.\ P \scong P_1 \mid P_2,\
     P_1 \sat \phi,\ P_2 \sat \psi\\[4pt]
  P \sat \alpha(\nu_1,\dots,\nu_n) &\text{iff}&
     \exists a_1,\dots,a_n.\ P \scong \alpha(a_1,\dots,a_n)
     \text{ and } a_i \sat \nu_i \text{ for all } i\\
  &&\text{for each shape } \alpha \neq \qq \text{ of arity } n\\[2pt]
  P \sat \qq(\nu,\psi) &\text{iff}&
     \exists a, S.\ P \scong \qq(a,S),\ a \sat \nu,\ S \sat \psi\\[4pt]
  P \sat \may{\phi}\psi &\text{iff}&
     \exists R \sat \phi.\ \exists P'.\
     P \mid R \red_{R} P' \text{ and } P' \sat \psi\\[4pt]
  P \sat_{\env} \fix X(\vec U).\,\phi
     &\text{iff}& P \in \bigcap \{\, F : \Phi_{\env}(F) \subseteq F \,\}\\
  P \sat_{\env} X(\vec\nu) &\text{iff}& P \in \env(X)(\sem{\vec\nu})
\end{array}
\]
where $\Phi_{\env}$ is the monotone operator on
$(\mathcal{P}(\mathsf{Name})^{n} \to \mathcal{P}(\mathsf{Proc}))$ induced by
$\phi$, and the least fixpoint exists by Knaster--Tarski since $X$ occurs
positively.
\end{definition}

\begin{definition}[Marked reaction]
\label{def:marked}
$P \mid R \red_{R} P'$ holds when there is a derivation of
$P \mid R \red P'$ by a single rule of Definition~\ref{def:red} in which
\emph{every} atom of $\comps{R}$ is a premise of that rule.
\end{definition}

\begin{remark}[Why the marking]
\label{rem:whymark}
Without it, $P \mid R \red P'$ is discharged by any internal step of $P$
whatsoever, and the formula $\phi$ describing the context does no work: every
term with a reduction would satisfy every $\may{\phi}\psi$. With it,
$\may{\nil}\psi$ is the internal-step modality --- $\comps{\nil}$ is empty, so
the condition is vacuous and the whole redex comes from $P$ --- and
$\may{\mm(\nu_1,\nu_2)}\psi$ says that $P$ offers a consumer at a name in
$\nu_1$ willing to accept a payload in $\nu_2$. The condition is on the
derivation and not on $P'$, which is where it has to be, since a residual
carries no record of which premises produced it. This follows
\cite[Def.~5.4]{softwhere}.
\end{remark}

\begin{remark}[One step, and only one]
$\red_{R}$ is a single rule application, not its reflexive transitive closure.
Over many steps a context could supply an atom, take a name from $P$, and act on
it, and the modality would be observing a protocol rather than a reaction. The
multi-step reading is available as a derived operator
(\S\ref{sec:derived}) and is the right thing to write when it is wanted.
\end{remark}

\begin{remark}[Bounded contexts, for free]
\label{rem:bdd}
Definition~\ref{def:marked} bounds $\comps{R}$ by the largest premise count of
any rule, which is four; and if $\comps{R}$ is nonempty, at least one premise
must come from $P$ unless $R$ alone forms a redex. This is what makes the modal
fragment tractable at all, and it is a gift of the combinator presentation: in a
calculus with prefixing, a reacting context can be arbitrarily large.
\end{remark}

\subsection{Positivity}

\begin{proposition}[Polarity]
\label{prop:polarity}
$\phi \mid \psi$, $\alpha(\dots)$, $\quo{(-)}$ and $\may{\phi}\psi$ are monotone
in every argument. The negative positions are exactly the scope of $\neg$ and
the antecedent of $\Rightarrow$.
\end{proposition}

\begin{proof}
Each clause of Definitions~\ref{def:satname} and~\ref{def:satproc} other than
$\neg$ and $\Rightarrow$ is a positive existential over the interpretations of
its arguments.
\end{proof}

Note that $\may{\phi}\psi$ is monotone in $\phi$ as well as in $\psi$: the
context is existentially quantified, so enlarging the class of admissible
contexts can only enlarge the satisfying set. Hence a fixpoint variable may
occur inside a modality's label, which \S\ref{sec:derived} uses.

%==============================================================================
\section{Derived forms}
\label{sec:derived}

\subsection{The spatial monoid}

\begin{proposition}[$\nil$ is the unit of $\mid$]
\label{prop:unit}
Up to logical equivalence, $(\Frm,\mid,\nil)$ is a commutative monoid:
$\phi \mid \nil \equiv \phi$, $\phi \mid \psi \equiv \psi \mid \phi$, and
$(\phi \mid \psi) \mid \chi \equiv \phi \mid (\psi \mid \chi)$. Moreover $\mid$
distributes over $\vee$ but not over $\wedge$, and
$\nil \equiv \neg(\neg\nil \mid \mathsf{true})$.
\end{proposition}

\begin{proof}
For the unit, $P \sat \phi \mid \nil$ iff $P \scong P_1 \mid P_2$ with
$P_1 \sat \phi$ and $P_2 \scong \nil$, iff $P \sat \phi$. Commutativity and
associativity descend from the abelian monoid structure of $\scong$.
Distribution over $\vee$ is immediate from the existential; failure over
$\wedge$ is the usual one, witnessed by
$\mm(a,v) \mid \mm(b,w)$ against
$(\mm(\quo{\mathsf{true}},\quo{\mathsf{true}}) \mid \mathsf{true}) \wedge
(\kk(\quo{\mathsf{true}}) \mid \mathsf{true})$.
The last identity says that $P$ is $\nil$ iff it has no nonempty component.
\end{proof}

\subsection{Counting}

\begin{definition}
\label{def:count}
For $k \geq 0$ and a formula $\theta$, write
\begin{align*}
  \#_{\theta} \geq k \ &\triangleq\ \underbrace{\theta \mid \cdots \mid
  \theta}_{k} \mid \mathsf{true},
  \qquad
  \#_{\theta} \leq k \ \triangleq\ \neg(\#_{\theta} \geq k{+}1),\\
  \#_{\theta} = k \ &\triangleq\ (\#_{\theta} \geq k) \wedge
  \neg(\#_{\theta} \geq k{+}1).
\end{align*}
Write $\shp$ for the disjunction of the structural productions of all shapes
other than $\mm$, each with every name argument $\quo{\mathsf{true}}$ and,
for $\qq$, payload $\mathsf{true}$; so $P \sat \shp$ iff $P$ is a single
non-message atom.
\end{definition}

So $\#_{\shp} \leq k$ is the atom-count measure of
\cite[Lem.~10.2]{rhocomb3}, written as a formula. The counting of
\cite{f1r3space}, which is primitive there, is here derived.

\subsection{Nominals, and the characteristic formula}

\begin{definition}[Characteristic formula]
\label{def:char}
For a closed term $T$ define $\chf{T}$ by structural recursion:
\[
\begin{array}{rcl}
  \chf{\nil} &=& \nil\\
  \chf{T_1 \mid T_2} &=& \chf{T_1} \mid \chf{T_2}\\
  \chf{\alpha(\quo{P_1},\dots,\quo{P_n})} &=&
     \alpha(\quo{\chf{P_1}},\dots,\quo{\chf{P_n}})
     \qquad (\alpha \neq \qq)\\
  \chf{\qq(\quo{P},S)} &=& \qq(\quo{\chf{P}},\chf{S}).
\end{array}
\]
\end{definition}

\begin{theorem}[Characteristic formulae are exact]
\label{thm:char}
For all closed $P$ and $T$: $P \sat \chf{T}$ iff $P \scong T$.
\end{theorem}

\begin{proof}
Induction on $|T|$. The base case is Definition~\ref{def:satproc}'s clause for
$\nil$. For an atom, $P \sat \alpha(\quo{\chf{P_1}},\dots)$ iff
$P \scong \alpha(a_1,\dots,a_n)$ with $a_i \sat \quo{\chf{P_i}}$, iff by
Definition~\ref{def:satname} $a_i \nequiv \quo{Q_i}$ with $Q_i \sat \chf{P_i}$,
iff by induction $Q_i \scong P_i$, iff $a_i \nequiv \quo{P_i}$ by
Proposition~\ref{prop:nequiv}. The case of $\qq$ is the same with the payload
compared at the process sort. For $T_1 \mid T_2$, satisfaction of
$\chf{T_1} \mid \chf{T_2}$ gives a split $P \scong P_1 \mid P_2$ with
$P_i \scong T_i$ by induction, and conversely the split witnessing
$P \scong T_1 \mid T_2$ satisfies the clause. Each recursive call is on a proper
subterm, so the induction is well founded.
\end{proof}

\begin{corollary}[Nominals]
\label{cor:nominal}
The process nominal of \cite[\S6]{softwhere} and the singleton namespace of
\cite{f1r3space} are both derived: $T$ as a formula is $\chf{T}$, and the
singleton $\{a\}$ for $a = \quo{P}$ is the namespace $\quo{\chf{P}}$.
\end{corollary}

\begin{corollary}[Maximal intensionality]
\label{cor:intensional}
$P \scong Q$ iff $P$ and $Q$ satisfy the same formulae. The logic is therefore
not adequate for any coarser equivalence, and in particular not for
$\bisim_{\enc{\Ctx_{\mathrm{rho}}}}$.
\end{corollary}

\subsection{Temporal operators}

\begin{definition}
\label{def:temporal}
\[
  \mst{\phi}\psi \ \triangleq\ \neg\may{\phi}\neg\psi,
  \qquad
  \Diamond\psi \ \triangleq\ \fix X().\ \psi \vee \may{\nil}X,
  \qquad
  \Box\psi \ \triangleq\ \neg\Diamond\neg\psi .
\]
\end{definition}

$\Diamond$ and $\Box$ range over internal reduction, since $\may{\nil}$ is the
internal step by Remark~\ref{rem:whymark}; $X$ occurs positively under
$\may{\nil}$ by Proposition~\ref{prop:polarity}, so the fixpoint is well formed.
The environment-facing analogue $\Diamond_{\phi}\psi \triangleq \fix X().\,\psi
\vee \may{\phi}X$ is available and is a genuinely different operator, since it
permits the context to be re-supplied at every step.

\begin{definition}[Bounded termination]
\label{def:term}
$\may{\nil}^{0}\phi \triangleq \phi$ and
$\may{\nil}^{j+1}\phi \triangleq \may{\nil}\may{\nil}^{j}\phi$; and
\[
  \mathsf{Halts}_{k} \ \triangleq\ \neg\may{\nil}^{k+1}\mathsf{true},
\]
which says no internal run from here has length exceeding $k$.
\end{definition}

%==============================================================================
\section{Model checking}
\label{sec:mc}

\subsection{The inhabitance relation}

The relation to be decided is $T \sat \phi$ for closed $T$ and closed $\phi$.
We call it \emph{inhabitance} when $\phi$ is read as a specification of a
position in the lattice and \emph{model checking} otherwise; they are the same
relation.

\begin{definition}[Hereditary material]
\label{def:sub}
$\Sub{T}$ is the least set of $\scong$-classes containing $T$ and closed under:
if $P \in \Sub{T}$ then every $Q$ with $\comps{Q} \subseteq \comps{P}$ is in
$\Sub{T}$; and if $\alpha(a_1,\dots,a_n) \in \comps{P}$ for some
$P \in \Sub{T}$, with $a_i = \quo{P_i}$, then each $P_i \in \Sub{T}$, and for
$\qq(a,S)$ also $S \in \Sub{T}$.
\end{definition}

\begin{lemma}
\label{lem:subfinite}
$\Sub{T}$ is finite, of size at most $2^{c} \cdot n$ where $n = |T|$ and $c$ is
the maximum number of top-level components of any hereditary subterm of $T$.
\end{lemma}

\begin{proof}
There are at most $n$ hereditary subterms, and the sub-multiset closure of a
term with $c$ components contributes at most $2^{c}$ classes. Descent is by the
subterm order, which is well founded on finite terms.
\end{proof}

\begin{theorem}[The modality-free fragment is decidable]
\label{thm:decstruct}
For closed $T$ and closed modality-free $\phi$, $T \sat \phi$ is decidable.
\end{theorem}

\begin{proof}
Every clause of Definition~\ref{def:satproc} other than the modality refers only
to members of $\Sub{T}$: the split clause to sub-multisets, the structural
clauses to quoted subterms and, for $\qq$, to the payload, and the name clause
to the quoted process. By Lemma~\ref{lem:subfinite} the carrier is finite, so
the operator $\Phi_{\env}$ of Definition~\ref{def:satproc} acts on the finite
lattice $\mathcal{P}(\Sub{T})^{\,\Sub{T}^{\,n}}$ and its least fixpoint is
reached by iteration in at most $|\Sub{T}|^{n+1}$ steps. Namespace arguments
range over $\{\quo{P} : P \in \Sub{T}\}$, again finite. Evaluate bottom up.
\end{proof}

\begin{remark}[Cost]
\label{rem:cost}
The bound of Lemma~\ref{lem:subfinite} is exponential in the number of parallel
components, and the split clause is where that is spent: deciding
$P \sat \phi \mid \psi$ naively enumerates the $2^{|\comps{P}|}$ splits. This is
the familiar cost of a spatial composition connective and is not an artefact of
this presentation. \S\ref{sec:occ} shows that the formulae the lattice actually
needs avoid it entirely.
\end{remark}

\subsection{The modal fragment}

\begin{proposition}[Decidability on the $\ev$-free sublattice, modulo witnesses]
\label{prop:decmodal}
Let $T \sat \Occ(\Atm)$ for some rule-closed $\Atm$ with $\ev \notin \Atm$.
Then the reachable set of $T$ is finite, and $T \sat \phi$ is decidable for
every closed $\phi$ whose modalities carry only labels $\phi_{\ell}$ with
$\sem{\phi_{\ell}}$ finite up to $\scong$.
\end{proposition}

\begin{proof}
By \cite[Prop.~10.7]{rhocomb3} every reduction sequence from $T$ has length at
most $|\comps{T}|_{\shp}$, and by Definition~\ref{def:red} the branching at each
step is bounded by the number of pairs and triples of components, so
$\mathrm{Reach}(T)$ is finite. Extend $\Sub{-}$ over it; the argument of
Theorem~\ref{thm:decstruct} then applies, with the modality clause evaluated by
enumerating the finitely many $R \sat \phi_{\ell}$ and, for each, the finitely
many marked derivations, which are bounded by Remark~\ref{rem:bdd}.
\end{proof}

The hypothesis on labels is the gap. A label such as
$\mm(\quo{\chf{P}},\quo{\mathsf{true}})$ constrains the subject to one name and
leaves the payload free, so $\sem{\phi_{\ell}}$ is infinite. The route to
removing the hypothesis is a depth-bounded distinguishability lemma:

\begin{conjecture}[Finite witnesses]
\label{conj:witness}
Define the quote depth $\mathrm{qd}(\phi)$ of a $\fix$-free formula as the
maximum nesting of $\quo{(-)}$ in it. Then for $\fix$-free $\phi_{\ell}$ and
$\psi$, satisfaction of $\may{\phi_{\ell}}\psi$ at $P$ is unchanged if $R$ is
restricted to range over a set of representatives of $\sem{\phi_{\ell}}$ of
quote depth at most $\mathrm{qd}(\phi_{\ell}) + \mathrm{qd}(\psi) + |P|$, which
is finite up to $\scong$ once the names available are those of $\Sub{P}$.
\end{conjecture}

The restriction to $\fix$-free formulae is not cosmetic: unfolding a fixpoint
raises quote depth without bound, so the bound must be re-derived per unfolding
or the conjecture restated as a bisimulation-up-to-depth argument. This is
listed as open item~\ref{open:witness}.

\begin{proposition}[Undecidability above]
\label{prop:undec}
With $\ev$ and $\dd$ present the calculus is Turing complete
\cite[\S7]{rhocomb3}, and $\Diamond \chf{T'}$ expresses reachability of $T'$.
So $T \sat \phi$ is undecidable for the full logic on the top vertex of the
computation axis.
\end{proposition}

\begin{remark}[The decidability frontier is the destructor edge]
\label{rem:frontier}
Propositions~\ref{prop:decmodal} and~\ref{prop:undec} say that the logic is
decidable exactly on the $\ev$-free half of the computation axis, which is the
edge \cite[Prop.~10.7]{rhocomb3} proves strict by strong normalisation. The
occupancy formula for that edge is therefore also the certificate that model
checking terminates. Should Conjecture~10.9 of \cite{rhocomb3} hold --- that
$\dd$-freedom gives strong normalisation even with $\ev$ --- the decidable
region extends to that vertex too, so settling it settles where model checking
stops working as well as where divergence starts.
\end{remark}

%==============================================================================
\section{Lattice inhabitance}
\label{sec:occ}

\subsection{Occupancy}

\begin{definition}[Hereditary atom set]
$\atoms{P}$ is the set of shapes $\alpha$ such that some atom of shape $\alpha$
occurs in $P$ at some quote depth: the least set with
$\alpha \in \atoms{P}$ whenever $\alpha(\vec a\,) \in \comps{P}$, and
$\atoms{Q} \subseteq \atoms{P}$ whenever $\quo{Q}$ is an argument of such an
atom or $Q$ is the payload of a $\qq$ in it.
\end{definition}

\begin{definition}[Occupancy formula]
\label{def:occ}
For $\Atm \subseteq \Atm_{\mathrm{full}}$,
\[
  \Occ(\Atm) \ \triangleq\ \fix X().\ \ \nil \ \vee\ (X \mid X)
  \ \vee \bigvee_{\alpha \in \Atm \setminus \{\qq\}}
     \alpha(\underbrace{\quo{X},\dots,\quo{X}}_{\mathrm{ar}(\alpha)})
  \ \vee\ [\,\qq(\quo{X},X) \ \text{if } \qq \in \Atm\,].
\]
\end{definition}

The formula is the atom list of Definition~\ref{def:syntax} with the disjuncts
for shapes outside $\Atm$ deleted, plus the two spatial disjuncts. The $\qq$
disjunct is the only one whose recursive occurrence appears at both sorts,
which is Remark~\ref{rem:migrate} again.

\begin{proposition}[Occupancy is hereditary containment]
\label{prop:occ}
$P \sat \Occ(\Atm)$ iff $\atoms{P} \subseteq \Atm$; equivalently, iff $P$ is a
term of the sub-calculus whose atom set is $\Atm$.
\end{proposition}

\begin{proof}
Right to left, the set of terms with $\atoms{P} \subseteq \Atm$ is closed under
the operator of Definition~\ref{def:occ}, hence contains its least fixpoint by
Knaster--Tarski, and conversely each disjunct preserves the property. Left to
right, observe that each recursive occurrence of $X$ is under a structural
production or under $\mid$, and that the argument of a structural production is
a proper subterm; so the unfolding descends the subterm order on the finite term
$P$ and terminates, at which point every atom encountered has been matched
against a disjunct of $\Atm$.
\end{proof}

\begin{corollary}
\label{cor:munu}
For $\Occ(\Atm)$ the least and greatest fixpoints coincide.
\end{corollary}

\begin{proof}
The unfolding descends the subterm order, which is well founded on finite
terms, so no infinite unfolding is consistent and $\nu$ adds nothing to $\mu$.
\end{proof}

\begin{theorem}[Occupancy is linear]
\label{thm:occlinear}
$T \sat \Occ(\Atm)$ is decidable in time linear in $|T|$.
\end{theorem}

\begin{proof}
By Proposition~\ref{prop:occ} the test is $\atoms{T} \subseteq \Atm$, which is
one traversal of $T$ descending into quotations. The split clause of
Remark~\ref{rem:cost} costs nothing here because the disjunct $X \mid X$ makes
$\Occ(\Atm)$ a homomorphic predicate: $P \mid Q \sat \Occ(\Atm)$ iff both
conjuncts do, so no enumeration of splits is required.
\end{proof}

Theorem~\ref{thm:occlinear} is the requirement of \S\ref{sec:what} discharged.
The occupancy formulae live in the modality-free fragment, so
Theorem~\ref{thm:decstruct} already gives decidability; what
Theorem~\ref{thm:occlinear} adds is that they land in the part of that fragment
which avoids the exponential entirely.

\subsection{Invariance, and a constraint on the lattice as drawn}

\begin{definition}[Rule-closed]
\label{def:ruleclosed}
$\Atm$ is \emph{rule-closed} if
\begin{enumerate}[label=(R\arabic*),leftmargin=2.6em,itemsep=2pt,topsep=3pt]
\item $\mm \in \Atm$;
\item $\fw \in \Atm$ whenever $\Atm$ meets $\{\bl,\br,\sy\}$;
\item $\consm \in \Atm \Rightarrow \mm \in \Atm$,
      $\consd \in \Atm \Rightarrow \dd \in \Atm$, and
      $\conss \in \Atm \Rightarrow \sy \in \Atm$.
\end{enumerate}
\end{definition}

\begin{proposition}[Invariance of occupancy]
\label{prop:subjred}
Let $\Atm$ be rule-closed. If $P \sat \Occ(\Atm)$ and $P \red P'$ then
$P' \sat \Occ(\Atm)$. Equivalently, $\Occ(\Atm) \Rightarrow \Box\,\Occ(\Atm)$
is valid.
\end{proposition}

\begin{proof}
By cases on Definition~\ref{def:red}, using Proposition~\ref{prop:occ}. The
router rules produce $\mm$ atoms and, for $\bl$, $\br$, $\sy$, an $\fw$ atom;
(R1) and (R2) cover these. The rule for $\ev$ produces the payload $P$ of a
message, whose atoms were already hereditarily present. The rule for $\qq$
produces $\mm(b,\quo{P})$, with $P$ the stored payload, again already present.
A constructor rule for shape $A$ produces $\mm(c,\quo{A(\vec v\,)})$, whose
atoms are those of $\vec v$ together with $A$ itself; (R1) and (R3) cover the
four cases, with $\consp$ needing nothing since parallel composition is not an
atom. Structural congruence introduces no atoms.
\end{proof}

\begin{proposition}[The two axes are not independent for occupancy]
\label{prop:notindep}
The lattice of \cite[Fig.~2]{rhocomb3} is drawn as a product of a computation
axis over $\{\dd,\ev\}$ with a constructor axis over the four constructors. At
the eight points carrying $\consd$ without $\dd$, the atom set is not
rule-closed, and occupancy is not invariant.
\end{proposition}

\begin{proof}
Take $\Atm$ with $\consd \in \Atm$, $\dd \notin \Atm$, and let
\[
  T \ =\ \consd(a,b,c,e) \mid \mm(a,u) \mid \mm(b,v) \mid \mm(c,w)
  \ \red\ \mm(e,\quo{\dd(u,v,w)}) \ =\ T'.
\]
Then $\atoms{T} \subseteq \Atm$ but $\dd \in \atoms{T'}$, so
$T \sat \Occ(\Atm)$ and $T' \not\sat \Occ(\Atm)$. The eight points are the two
computation points omitting $\dd$ crossed with the four constructor points
including $\consd$.
\end{proof}

\begin{remark}[Three ways out]
\label{rem:waysout}
Proposition~\ref{prop:notindep} is a statement about the sub-calculus, not about
the logic, and it is the logic's first return on the investment. A sub-calculus
that can build the name $\quo{\dd(u,v,w)}$ but has no $\dd$ atom is not closed
under its own reduction, because the names of a sub-calculus are quotations of
its own processes. Three repairs are available and the note does not choose
between them: (i) keep $\dd$ in the base for those points, so the constructor
axis is a product only over the remaining three; (ii) read occupancy
non-hereditarily, at the cost of Proposition~\ref{prop:occ} no longer saying
``is a term of''; or (iii) accept the eight points as expressively meaningful
but not reduction-closed, and state the product claim for the expressive
ordering only. Option (i) matches the intent of \cite[\S10.4]{rhocomb3}, where
$\dd$ is discussed as a source of copies rather than as a shape a name may
carry.
\end{remark}

\subsection{Name closure: the constructor axis}

Occupancy says which atoms may appear. The constructor axis wants something
sharper: which \emph{names} can ever appear, which is
\cite[Lem.~10.1]{rhocomb3}. That lemma is a namespace fixpoint, and its
disjuncts are in bijection with the constructors present.

\begin{definition}[Closure namespace]
\label{def:cl}
Let $T$ be closed with $\nms{T} = \{\quo{Q_1},\dots,\quo{Q_r}\}$ and let
$S \subseteq \{\consp,\consm,\consd,\conss\}$. Put
\begin{align*}
  \Cl_{S}(T) \ \triangleq\ \quo{\big(\fix X().\ \kappa_{S}\big)},
  \qquad \text{where } \kappa_{S} \ =\ \
  & \bigvee_{i=1}^{r} \chf{Q_i}\\
  \vee\ & [\,X \mid X & &\text{if } \consp \in S\,]\\
  \vee\ & [\,\mm(\quo{X},\quo{X}) & &\text{if } \consm \in S\,]\\
  \vee\ & [\,\dd(\quo{X},\quo{X},\quo{X}) & &\text{if } \consd \in S\,]\\
  \vee\ & [\,\sy(\quo{X},\quo{X},\quo{X}) & &\text{if } \conss \in S\,].
\end{align*}
\end{definition}

\begin{definition}[Confinement]
\label{def:conf}
\[
  \Conf(U) \ \triangleq\ \fix X(U).\ \ \nil \ \vee\ (X(U) \mid X(U))
  \ \vee \bigvee_{\alpha \in \Atm_{\mathrm{full}} \setminus \{\qq\}}
    \alpha(\underbrace{\theta_U,\dots,\theta_U}_{\mathrm{ar}(\alpha)})
  \ \vee\ \qq(\theta_U, X(U)),
\]
where $\theta_U \triangleq U \wedge \quo{X(U)}$ --- a namespace, well formed by
Proposition~\ref{prop:namecollapse}.
\end{definition}

$\Conf(U)$ says that every name occurring at any depth lies in $U$, and that the
processes those names quote are themselves confined. This is the one place the
parameterised fixpoint earns its keep: the parameter $U$ must be carried
unchanged through the recursion, and a nullary fixpoint cannot do it, because
the namespace being tested is not definable from the term at hand.

\begin{proposition}[Name closure, in the logic]
\label{prop:nameclosure}
Let $T$ be closed, $S$ the set of constructors in $\Atm$, and $\Atm$ rule-closed.
Then $T \sat \Conf(\Cl_{S}(T))$, and
$\Conf(\Cl_{S}(T)) \Rightarrow \Box\,\Conf(\Cl_{S}(T))$ is valid at $T$.
\end{proposition}

\begin{proof}
The first claim is immediate: every name of $T$ is some $\quo{Q_i}$, matched by
the corresponding disjunct of Definition~\ref{def:cl} via
Theorem~\ref{thm:char}. The second is \cite[Lem.~10.1]{rhocomb3} transcribed.
Every rule other than a constructor rule has on its right only names occurring
on its left; a constructor rule for shape $A$ emits
$\mm(c,\quo{A(\vec v\,)})$ with $\vec v$ already present, and $\quo{A(\vec v\,)}$
lies in $\Cl_{S}(T)$ precisely because the disjunct for $\cons_{A}$ is present
in Definition~\ref{def:cl} exactly when $\cons_{A} \in S$.
\end{proof}

\begin{corollary}[Root exclusion]
\label{cor:root}
Write $\mathrm{Root}_{A}$ for the formula
$\neg\nil \mid \neg\nil$ when $A$ is parallel composition, and
$A(\quo{\mathsf{true}},\dots,\quo{\mathsf{true}})$ otherwise. If
$\cons_{A} \notin S$ and no name of $\nms{T}$ satisfies
$\quo{\mathrm{Root}_{A}}$, then
\[
  T \ \sat\ \Box\,\neg\big(\ \mm(\quo{\mathrm{Root}_{A}},\quo{\mathsf{true}})
  \mid \mathsf{true}\ \big),
\]
that is, no message on a subject of root $A$ is ever emitted.
\end{corollary}

\begin{proof}
By Proposition~\ref{prop:nameclosure} every name reachable from $T$ lies in
$\Cl_{S}(T)$, whose defining disjunction has no clause with root $A$ and whose
base disjuncts are excluded by hypothesis. A message with such a subject would
require one.
\end{proof}

Corollary~\ref{cor:root} is the converse half of
\cite[Prop.~10.5]{rhocomb3} in logical form. The forward half --- that the
behaviour \emph{is} realised when $\cons_{A}$ is present --- is a modal formula:

\begin{definition}[Separation witness]
\label{def:witness}
Let $A$ be a node shape of arity $n$, and let the names
$a$, $v$ and $k_2,\dots,k_n$ be fixed. Put
\[
  W_{A} \ \triangleq\ \may{\ \mm(\quo{\chf{P_a}},\quo{\chf{P_v}})\ }\
  \Diamond\Big(\ \mm\big(\quo{\chf{A(v,k_2,\dots,k_n)}},
  \quo{\mathsf{true}}\big) \mid \mathsf{true}\ \Big),
\]
where $a = \quo{P_a}$ and $v = \quo{P_v}$.
\end{definition}

$W_{A}$ is $\mathrm{Spec}_{A}$ of \cite[Def.~10.4]{rhocomb3} written as a
formula: on receiving $v$ at $a$, emit a message on the constructed subject. By
Corollary~\ref{cor:root}, $W_{A}$ is unsatisfiable under
$\Occ(\Atm) \wedge \Conf(\Cl_{S}(T))$ when $\cons_{A} \notin S$ and $v$ is
chosen outside $\nms{T}$, and it is satisfiable when $\cons_{A}$, $\dd$, $\br$
and $\mm$ are present. The choice of $v$ remains a side condition of the
separation argument rather than part of the formula, since it quantifies over
the term being separated.

\subsection{The formulae, point by point}

\begin{definition}[Base]
$\Atm_{0} \triangleq \{\mm,\kk,\fw,\bl,\br,\sy,\qq\}$: the medium, the store,
and the five routers that are not $\dd$.
\end{definition}

\begin{definition}[Lattice points]
\label{def:points}
For $c \subseteq \{\dd,\ev\}$ and $S \subseteq \{\consp,\consm,\consd,\conss\}$
with $S \neq \emptyset \Rightarrow \consp \in S$, put
$\Atm_{c,S} \triangleq \Atm_{0} \cup c \cup S$, and
\[
  \mathsf{At}_{c,S} \ \triangleq\ \Occ(\Atm_{c,S}),
  \qquad
  \mathsf{In}_{c,S}(T) \ \triangleq\
  \Occ(\Atm_{c,S}) \wedge \Conf(\Cl_{S}(T)) .
\]
\end{definition}

$\mathsf{At}_{c,S}$ is the occupancy formula: it holds of exactly the terms of
the sub-calculus. $\mathsf{In}_{c,S}(T)$ adds the name-closure invariant, which
is what the strictness arguments consume; it is term-relative because the
closure namespace is generated from $\nms{T}$.

\begin{center}
\begin{tabular}{@{}llll@{}}
\toprule
axis & point & formula & by\\
\midrule
computation & $\ev \notin c$ & $\Occ(\Atm_{c,S})$, $\ev$ disjunct deleted
  & Prop.~\ref{prop:occ}\\
 & with $\#_{\shp} \leq k$ & $\Rightarrow \mathsf{Halts}_{k}$
  & \cite[Prop.~10.7]{rhocomb3}\\
 & $\dd \notin c$ & $\Occ(\Atm_{c,S})$, $\dd$ disjunct deleted
  & Prop.~\ref{prop:occ}\\
\midrule
constructors & $\cons_{A} \notin S$ & $\Occ(\Atm_{c,S})$, $\cons_{A}$ disjunct
  deleted & Prop.~\ref{prop:occ}\\
 & name closure & $\Conf(\Cl_{S}(T))$, $A$ disjunct deleted
  & Prop.~\ref{prop:nameclosure}\\
 & separation & $\neg W_{A}$ & Cor.~\ref{cor:root}\\
\midrule
both & invariance & $\Box$ of the above, $\Atm$ rule-closed
  & Prop.~\ref{prop:subjred}\\
 & the eight bad points & $\consd \in S,\ \dd \notin c$: not invariant
  & Prop.~\ref{prop:notindep}\\
\bottomrule
\end{tabular}
\end{center}

\begin{proposition}[The computation axis, in the logic]
\label{prop:compaxis}
$\Occ(\Atm_{c,S}) \wedge (\#_{\shp} \leq k) \Rightarrow \mathsf{Halts}_{k}$ is
valid whenever $\ev \notin c$.
\end{proposition}

\begin{proof}
\cite[Lem.~10.2]{rhocomb3}: every rule but that for $\ev$ strictly decreases the
number of non-message atoms, and $\#_{\shp}$ counts exactly those by
Definition~\ref{def:count}. Hence no internal run exceeds $k$ steps, which is
Definition~\ref{def:term}.
\end{proof}

\begin{remark}[What the formulae do not decide]
\label{rem:notdecide}
$\mathsf{At}_{c,S}$ decides whether $T$ \emph{is} a term of the sub-calculus.
It does not decide whether $T$ is \emph{equivalent to} one, which is the
lattice's own ordering. The gap between the two is exactly the content of the
strictness results: Proposition~\ref{prop:compaxis} and
Corollary~\ref{cor:root} are the bridges, and they run one way, ruling out
behaviours below a point rather than characterising the terms at it. A
characterisation of the semantic position --- ``$T$ is $\bisim$-equivalent to
some term satisfying $\mathsf{At}_{c,S}$'' --- is not expressible in this logic
and is very unlikely to be decidable; see open item~\ref{open:semantic}.
\end{remark}

%==============================================================================
\section{What the logic does not do}
\label{sec:notdo}

\paragraph{The observational lattice.} \cite[Fig.~1]{rhocomb3} is a second
lattice over the same subsets, on which $\qq$ and $\mm$ occupy different points
and which collapses against encoded observers. Its edges are receptivity
statements: $\mm(a,v)$ is consumable by every consumer at $a$, while $\qq(a,P)$
is consumable only by $\fw$ \cite[Rem.~3.6]{rhocomb3}. Receptivity is the
$\forall\exists$ modality $\rcp{K}\phi$ of \cite[\S7]{softwhere}, and it is not
derivable from $\may{\phi}\psi$ and negation, which give only the
$\exists\exists$ and $\forall\forall$ readings --- may and must quantify over
the same player. The logic as specified therefore sees the $\qq$--$\mm$
distinction \emph{structurally}, through the separate productions
$\qq(\nu,\psi)$ and $\mm(\nu_1,\nu_2)$, and not \emph{behaviourally}. For
occupancy that is enough. For anything about the observational ordering it is
not, and the repair is to add $\rcp{\phi}\psi$ as a second primitive.

\paragraph{Adequacy.} By Corollary~\ref{cor:intensional} the logic separates
terms up to $\scong$, so it is not a characteristic logic for
$\bisim_{\enc{\Ctx_{\mathrm{rho}}}}$, nor for any coarser relation. This is a
design choice and not an oversight: an occupancy formula must see the shape of
a term, and any logic that does will be at least as discriminating as its
structural layer. Identifying the extensional sublogic --- the fragment
invariant under $\bisim_{\Obs}$ --- is the natural companion question and is
open item~\ref{open:extensional}. The same tension in its general form is the
subject of \cite{extnote}.

\paragraph{Minimal labels.} The labels of $\may{\phi}\psi$ are enabling
contexts, following \cite[\S6]{rhocomb3}; the minimal-label discipline via
relative pushouts is \cite{bihol}'s and is not used here. Nothing in this note
needs it, and Remark~\ref{rem:bdd} means the enabling contexts are already
small.

\paragraph{Cost.} The logic has no grading. The name-growth question of
\cite[\S13]{rhocomb3} --- that with constructors the names are no longer
subterms of the program --- shows up here as the observation that $\Cl_{S}(T)$
is an infinite namespace with a finite presentation, which is exactly the
situation a graded clause is built for. That is the bridge to
\cite{softwhere}'s stage two and is not pursued.

%==============================================================================
\section{Open items}
\label{sec:open}

\begin{enumerate}[leftmargin=1.8em,itemsep=3pt]

\item \label{open:witness} \textbf{Finite witnesses for the modality.}
Conjecture~\ref{conj:witness}, and its extension past the $\fix$-free fragment.
Without it Proposition~\ref{prop:decmodal} carries a hypothesis on labels that
excludes the labels one actually writes.

\item \label{open:eight} \textbf{The eight non-invariant points.}
Proposition~\ref{prop:notindep} and Remark~\ref{rem:waysout}. This is a question
for \cite{rhocomb3} rather than for the logic, and the answer changes
\cite[Fig.~2]{rhocomb3}.

\item \label{open:extensional} \textbf{The extensional sublogic.} Which formulae
are invariant under $\bisim_{\enc{\Ctx_{\mathrm{rho}}}}$? The modality-only
fragment is the obvious candidate and is not obviously closed under the
fixpoint.

\item \label{open:semantic} \textbf{Semantic position.} Is there any useful
approximation to ``equivalent to a term of the sub-calculus'' that the logic can
express? Remark~\ref{rem:notdecide}.

\item \label{open:recept} \textbf{Receptivity.} Adding $\rcp{\phi}\psi$, and
whether Theorem~\ref{thm:decstruct} survives it. The $\forall$ over contexts is
the same infinity that Decision~\ref{dec:omit} removed with the guarantee, so
the answer is probably that it survives only under
Conjecture~\ref{conj:witness}.

\item \label{open:impl} \textbf{Implementation.} The structural fragment is a
bottom-up tree-automaton run and should fall out of the existing
\texttt{f1r3space} crate \cite{f1r3space} with the restriction connectives
deleted and the atom productions added. Theorem~\ref{thm:occlinear} says the
occupancy checker is a separate and much simpler artefact, and it is the one the
lattice needs.

\item \label{open:conj} \textbf{Duplication.} \cite[Conj.~10.9]{rhocomb3}.
Remark~\ref{rem:frontier} gives a second reason to want it settled.

\end{enumerate}

%==============================================================================
\sloppy
\begin{thebibliography}{99}

\bibitem{bihol} C.\,B. Wells, M. Stay and L.\,G. Meredith.
\emph{Behavior in higher-order languages: admissible observers, relative
pushouts, and the encoding of extended theories.}
In preparation, F1R3FLY.io, 2026.

\bibitem{Caires04} L. Caires.
\emph{Behavioral and spatial observations in a logic for the $\pi$-calculus.}
FOSSACS 2004, LNCS 2987, 72--89.

\bibitem{CairesCardelli03} L. Caires and L. Cardelli.
\emph{A spatial logic for concurrency (part I).}
Information and Computation 186(2):194--235, 2003.

\bibitem{extnote} L.\,G. Meredith.
\emph{Every constructor an interaction: observer extensions, and the
discriminating power of spatial formulae in interactive graph-structured lambda
theories.}
F1R3FLY.io, 2026.

\bibitem{f1r3space} F1R3FLY.io.
\emph{F1R3Space --- semantics of the rho model checker.}
Draft 0.1, \texttt{F1R3FLY-io/\allowbreak F1R3Space}, \texttt{MC-Rholang/}, 2026.

\bibitem{MeredithRadestock05} L.\,G. Meredith and M. Radestock.
\emph{A reflective higher-order calculus.}
Electronic Notes in Theoretical Computer Science 141(5):49--67, 2005.

\bibitem{rhocomb3} L.\,G. Meredith and M. Stay.
\emph{Name-free combinators for the rho calculus: a repair, a linear encoding,
and the one operation that was missing.}
Draft 3, \texttt{F1R3FLY-io/\allowbreak publications}, \texttt{rho-combinators/}, 2026.

\bibitem{softwhere} F1R3FLY.io.
\emph{Context-labelled hybrid logic for a rho calculus with where-clauses.}
Specification for the F1R3Ware checker,
\texttt{F1R3FLY-io/\allowbreak F1R3Space}, \texttt{Softwhere/\allowbreak where-\allowbreak context-\allowbreak labels.tex},
2026.

\bibitem{Yoshida02} N. Yoshida.
\emph{Minimality and separation results on asynchronous mobile processes:
representability theorems by concurrent combinators.}
Theoretical Computer Science 274(1--2):231--276, 2002.

\end{thebibliography}

\end{document}
